\section{Thông tin đồ án}

\subsection{Thông tin thành viên}

\begin{table}[!ht]
\centering
\renewcommand{\arraystretch}{1.25}
\begin{tabular}{|>{\centering\arraybackslash}m{5cm}|>{\centering\arraybackslash}m{3cm}|>{\centering\arraybackslash}m{8cm}|}
    \hline
    \Large \textbf{Tên} & \Large \textbf{MSSV} & \Large \textbf{Email} \\
    \hline
    \Large Ngũ Kiệt Hùng & \Large 22127134 & \Large nkhung22@clc.fitus.edu.vn \\
    \hline
    \Large Võ Thành Nghĩa & \Large 22127295 & \Large vtnghia22@clc.fitus.edu.vn \\
    \hline
    \normalsize Phạm Nguyên Hải Long  & \Large 22127248 & \Large pnhlong22@clc.fitus.edu.vn \\
    \hline
\end{tabular}
\end{table}

\subsection{Giới thiệu đồ án}
\quad Mục đích của Vision Document là để thu thập, phân tích và xác định các yêu cầu cấp cao cùng những tính năng mong muốn cho một hệ thống giám sát giao thông sử dụng Edge AI. Hệ thống được thiết kế nhằm xây dựng một giải pháp giám sát tình trạng ùn tắc giao thông mà không sử dụng các dịch vụ của bên thứ 3 thông qua việc triển khai giải pháp phát hiện, phân loại và đếm số lượng phương tiện di chuyển qua một tuyến đường chính theo thời gian thực bằng  trên các phần cứng nhỏ với chi phí thấp .

\quad Tài liệu này ghi nhận các quan điểm và nhu cầu của các bên liên quan như quản lý giao thông đô thị và các nhà phát triển, nhằm cung cấp một cái nhìn rõ ràng về mục tiêu hệ thống và tính ứng dụng thực tế của nó trong các thành phố thông minh.

\quad Tài liệu sẽ bắt đầu với phần Định hướng, được trình bày tại Mục \nameref{sec:vision} nơi nhóm phát triển sẽ hình thành vấn đề dựa trên yêu cầu của người dùng và các bên liên quan dưới dạng các tác vụ cụ thể. Tiếp theo, tài liệu sẽ cung cấp cái nhìn tổng quan về các bên liên quan và người dùng trong Mục \nameref{user_description} bao gồm các yêu cầu và ràng buộc từ phía họ.

\quad Mục \nameref{product_overview} sẽ trình bày khái quát về ứng dụng dịch vụ và mối quan hệ với các dịch vụ bên thứ ba, cũng như các giả định về sự phụ thuộc trong kiến trúc hệ thống. Mục \nameref{sec:product_features} sẽ mô tả chi tiết hơn về các tính năng cốt lõi được kỳ vọng trong phiên bản triển khai thực tế của dịch vụ.
Cuối cùng, Mục \nameref{feature_request} sẽ xác định những yêu cầu về mặt vận hành của dịch vụ, các ràng buộc về thời gian hoạt động (uptime), cũng như hạ tầng vận hành được khuyến nghị cho người dùng và các yêu cầu dành cho người dùng cuối.

\subsection{Những thay đổi so với phiên bản gần nhất}
\begin{table}[H]
    \centering
    \caption{So sánh phiên bản cũ và phiên bản hiện tại}
    \begin{tabular}{|p{8cm}|p{4cm}|p{4cm}|}
    \hline
    \textbf{Mục} & \textbf{Phiên bản cũ} & \textbf{Phiên bản hiện tại} \\ \hline
    Thay đổi công cụ & Sử dụng Google Docs & Sử dụng Latex \\ \hline
    Mục tiêu 1: Chi tiết hoá Measurable (Đo lường được) & Không có F1 & Ghi rõ số liệu của Precision và F1 \\ \hline
    Mục tiêu 2: Thay đổi mục Specific & Người dùng có thể tương tác & Người dùng chỉ có thể quan sát \\ \hline
    Mục tiêu 3: Giảm độ chính xác ghi rõ là F1 &  Độ chính xác >= 70\%–80\%  & Độ chính xác (F1) >= 60\%  \\ \hline
    Thay đổi sơ đồ kiến trúc hệ thống (theo yêu cầu của thầy, vì ảnh trước các edges hơi nhỏ) & Sơ đồ có nhiều edge & Sơ đồ có 1 edge duy nhất \\ \hline
    Quy trình CICD & Không có & Thêm ảnh và mô tả về quá trình CICD \\ \hline
    \end{tabular}
\end{table}





\pagebreak